\section{Chunk Aggregation}
\label{app:aggregation}

This appendix measures the aggregator of Section~\ref{subsec:reanchoring}
against submitting the chunk proofs one by one. The measurements were taken
offline on a 10-core, 32\,GiB macOS workstation, not on Testbed~A, with
nargo 1.0.0-beta.22 and \texttt{bb} 5.0.0-nightly.20260522. The chain is
synthetic: it has one deterministic header sequence with real ed25519
precommits and one validator set, so that consecutive chunks chain through
their roots and heights. Aggregate proving and verification are averaged over
five trials per interval length; each leaf proof was timed once. Absolute times
are not comparable to Table~\ref{tab:reanchor-scaling}, which was measured on a
different host.

\begin{table}[h]
\caption{Chunk Aggregation vs.\ Submitting Chunk Proofs One by One ($K{=}4$)}
\label{tab:aggregation}
\centering
\footnotesize
\begin{tabular}{@{}rrrrrrr@{}}
\toprule
 & \multicolumn{3}{c}{\textbf{One by one}} & \multicolumn{3}{c}{\textbf{Aggregate}} \\
\cmidrule(lr){2-4}\cmidrule(lr){5-7}
$L$ & Prove & Bytes & Verify & Prove & Bytes & Verify \\
\midrule
256  & 31.5 s  & 14{,}656 & 19.8 ms & 57.2 s  & 14{,}656 & 20.2 ms \\
512  & 60.5 s  & 29{,}312 & 39.5 ms & 89.9 s  & 14{,}656 & 20.2 ms \\
1024 & 120.1 s & 58{,}624 & 79.1 ms & 146.7 s & 14{,}656 & 20.2 ms \\
\bottomrule
\end{tabular}
\\[2pt]
\footnotesize Verify is the validators' embedded WebAssembly verifier, measured
per proof (30 iterations, three runs); the one-by-one column is $n$ times the
leaf cost. Aggregate prove includes the leaf proofs. At $L{=}256$ a single
leaf proof is submitted and the aggregate row is shown for the overhead only.
\end{table}

The aggregator verifies the chunk proofs in-circuit and checks that chunk
$i$ starts at the root and height where chunk $i{-}1$ ended, under the same
validator-set hash. It exposes the same public inputs as a single chunk, so an
aggregate proof stands in for a leaf proof, and its capacity $K{=}4$ is fixed
at compile time. Unbounded incremental recursion would require the step circuit
to receive its own verification-key hash as a public input, which we did not
implement. Recursion is applied at the chunk level because in-circuit
verification of one UltraHonk proof costs about 0.7 million gates, whereas
hashing one header costs about 12 opcodes: recursion over per-header proofs
measured 650--1{,}200$\times$ slower than direct hashing at matched sizes.

Table~\ref{tab:aggregation} shows that submitted bytes and verification time of
the aggregate do not depend on $L$ up to $K \cdot 256 = 1024$, while the
one-by-one cost grows linearly. The price is one extra proof: the aggregator
circuit has 3{,}002{,}299 gates and proves in 25.7--29.5\,s using about 5.2\,GB,
against 3{,}234{,}226 gates, 30.0--31.5\,s, and 5.7--5.9\,GB for a leaf. Total
proving time rises by 49\% at $L{=}512$ and 22\% at $L{=}1024$. Chunk proofs
were produced sequentially; proving them in parallel was not measured. Every
chunk boundary height needs a round-zero commit, so a round skip at an interior
boundary prevents building that chunk's witness.

The aggregator rejects each of eleven malformed witnesses: reversed or repeated
chunks, a wrong total count, final root, first height, or validator-set hash,
a height gap between chunks, a different validator set in the second chunk,
an extension through padding slots, and a chunk count outside $[1,K]$. The
valid two-chunk witness is accepted, and a 512-header aggregate passes every
\texttt{SubmitRecoveryProof} binding against fixture chain state. On the Docker
testbed, aggregate proofs of three and two chunks (514 and 512 headers) were
accepted by \texttt{SubmitRecoveryProof} on a live chain; whether the chain then
leaves \textsc{Recovering} depends on the unproven-tail bound, as the
Discussion reports.

\FloatBarrier
